% !TEX root = ./main.tex

We describe two implementations of the Filtered Stepper Calculus: a
compact reference implementation in OCaml that directly mirrors the
formal rules, and the full integration into the Hazel live programming
environment.

\subsection{Reference Implementation}\label{sec:impl-mini}

To validate the formal semantics presented in Section~\ref{sec:filter},
we built a standalone reference stepper in OCaml. It implements a
substitution-based small-step evaluator for the simply typed lambda
calculus extended with a fixpoint operator and the full filter calculus.
The core is roughly 560 lines of OCaml, structured as a direct
transliteration of the formal judgments: expressions and patterns are
algebraic data types, actions and gas are simple sums, and filter and
residue expressions are constructors of the expression type.

A single stepping function orchestrates the pipeline described in
Section~\ref{sec:dynamics}, calling instrumentation, optimization,
decomposition, action selection, transition, decay, and composition
in sequence. If action selection yields \(\ActSkip\), or if the redex
is a filter or residue elimination (the \(\Strippable{}\) case), the
stepper performs the transition silently and recurses. Otherwise it
returns the list of available redexes to the caller.

Because the reference implementation uses substitution-based evaluation
and named variables, it corresponds closely to the on-paper calculus and
is easy to audit against the formal rules. We test the stepper against
known results (e.g., verifying that \texttt{fac(5)} evaluates to 120).
On the common subset of the two calculi (lambda calculus with addition,
fixpoints, and filters), the reference implementation also serves as a
differential test oracle for the Hazel integration; any divergence
between the two indicates a bug in the production stepper. Hazel
features beyond that subset---pattern matching, sum and product types,
type ascriptions, typed holes---are covered by separate hand-written
unit tests against expected trace outputs.

\subsection{Hazel Integration}\label{sec:impl-hazel}

Hazel internally uses a big-step, environment-based evaluator for
performance: it avoids redundant substitutions and enables tail-call
optimization. However, because the course teaches substitution-based
equational semantics, the stepper must present traces as if
substitution were performed eagerly.

To keep the two evaluation styles consistent, the big-step evaluator
and the small-step stepper share a common transition relation: each
redex's reduction case is defined in one place, and the two
interpreters differ only in their driver loop (recursive
descent for the big-step evaluator versus single-redex extraction
for the stepper). The stepper works over Hazel's internal expression
representation, which includes closures (pairs of an expression and
an environment). During decomposition, closures are traversed like
ordinary expressions, with the environment carried in the evaluation
context. When the stepper displays a step to the user, a
post-processing pass replaces all bound variables with their
environment bindings, producing output that looks substitution-based.

The filter system in Hazel mirrors the reference implementation but
operates over Hazel's richer expression language, which includes
pattern matching, sum and product types, type ascriptions, and typed
holes. Filter matching uses a semantic equality check with special
accommodations: fixpoint wrappers and type ascriptions are ignored
during comparison.

Beyond user-written filters, Hazel exposes UI toggles that classify
certain step kinds (fixpoint unrolling, type-ascription elimination,
pattern-match redex selection) as skipped-by-default for the trace,
while a few others (closure-wrapping, residue bookkeeping) are
unconditionally hidden because they correspond to internal machinery
rather than user-visible reductions.

% We moved the content of performance section to be alongside with the
% Optimization rule (O-*).

\subsection{Performance}\label{sec:impl-perf}

We have already discussed the cost of instrumentation in
Section~\ref{sec:dynamics}. The instrumentation rules walk the whole
expression, and at each node test the active filter's pattern against that
node---the I-Y/I-N split. With $k$ filters in scope and an expression of size
$n$, this is $O(n \cdot k \cdot |p|)$ work per step, where $|p|$ is the largest
pattern. The $k$ factor is unavoidable, since a deeply nested filter must still
be considered at every node it covers.

The structural problem is what happens when several nested filters
match the same redex. Two effects compound. First, within a single
instrumentation pass, the recursive structure of \textsc{I-Filter}
re-walks already-instrumented sub-expressions and can stack residues
exponentially in the nesting depth (Section~\ref{sec:dynamics}).
Second, those residues are not transient: they persist as part of
the expression into the next step, where re-instrumentation walks
over them and may add more on top. Without intervention, residue
chains grow with the \emph{history} of stepping as well as with
filter nesting depth, so the cost of decomposition and action
selection drifts away from $O(n)$ and keeps climbing.

The optimization pass cuts this back. \textsc{O-Residue-Inner} and
\textsc{O-Residue-Outer} collapse each maximal chain of adjacent
residues to its highest-priority survivor, so residue depth per chain
stays at $O(1)$ and the total residue count stays at $O(n)$ regardless
of how many filters are in scope. Decomposition and action selection
then work on an expression tree of size $O(n)$ rather than
$O(n \cdot k)$. The $|p|$ factor in instrumentation remains the
asymptotic bottleneck per step. This is unsurprising, since pattern
matching is the work the user is actually paying for.

%%% Local Variables:
%%% mode: LaTeX
%%% TeX-master: "main"
%%% End:
